\documentclass[10pt]{article}
\usepackage[T1]{fontenc}
\usepackage{lmodern}
\usepackage[a4paper,margin=25mm]{geometry}
\usepackage{booktabs}
\usepackage{array}
\usepackage{microtype}
\usepackage[hidelinks]{hyperref}
\setlength{\parindent}{0pt}
\setlength{\parskip}{0.55em}

\title{A Complete Teacher Failure and Offline Repair Episode}
\author{GLEAM bridge example from the procurement reasoning environment}
\date{}

\begin{document}
\maketitle

\section{Purpose and evidence boundary}

This document reconstructs one complete reasoning episode for the following benchmark item.
\begin{center}
\texttt{L2::bridge\_join\_0583\#L2a}
\end{center}
It connects the benchmark question, the teacher call, the first deterministic execution, the
offline rejection, the repair messages, the accepted program, and its deterministic replay.

Three evidence statuses are kept separate.  Stored artifacts are copied from the historical JSONL
corpus.  Reconstructed prompts are regenerated from the committed prompt builder and the stored
case fields because the compact corpus did not retain every provider request and response byte for
this item.  Replay observations are recomputed with the stored accepted plan on the current
versioned graph.  No reconstructed text is represented as an original provider log.

\section{Question and evaluation target}

\begin{quote}
Just looking at the relevant procurement records, how many contract notices were published by
buyers that went on to award a contract to GLEAM SERVICE CENTRE LTD?
\end{quote}

The required answer type is a count.  The benchmark generator executes the latent bridge
specification against its graph snapshot and stores 1,623 as the evaluation oracle.  This value is
available to the offline evaluator but is not included in a planner or repair prompt.

The intended computation is
\begin{quote}\ttfamily
awards to GLEAM\par
$\rightarrow$ distinct buyers on those awards\par
$\rightarrow$ all procurement records published by those buyers\par
$\rightarrow$ count.
\end{quote}

\section{Teacher Stage 1}

The fixed Stage 1 system message is
\begin{quote}\small
You are the typed intent programmer for a UK public-procurement KGQA system. Convert the question
into a small computation program. Do not answer the question. Use only information explicitly
present in the question. The output shape is enforced by a strict JSON schema.
\end{quote}

The case specific user object contains the question, allowed procurement slots, operation
guidance, hard rules, and retrieved schema context.  Its bridge guidance states that a bridge must
use source records, a distinct entity set, a bound target record set, and a final aggregate or
selection.  The exact object is generated by
\texttt{question\_intent\_program\_messages} in
\path{src/procurement_graph/reasoning/typed_planning.py}.

The stored Stage 1 response declares a numeric count but proposes only one record filtering step:
\begin{verbatim}
answer_signature:
  operation: count
  value_type: number
  answer_field_text: ""

program:
  A = filter_records(
        supplier is awarded contract to = GLEAM SERVICE CENTRE LTD,
        buyer publishes contract notice)
  returns: record_set

answer_step: A
\end{verbatim}

The response also says that the schema context does not expose a distinct contract notice type or
field.  It preserves the supplier literal, but it does not represent the central relational
requirement.  GLEAM first determines a set of buyers, and those buyers determine the population to
count.

\section{Deterministic compilation and the rejected program}

The stored Stage 1 object satisfies the typed intent program contract.  It therefore enters the
deterministic intent compiler rather than requiring another language model to construct the initial
graph.  Recompiling the stored object with the committed compiler yields a graph exactly equal,
field for field, to the historical rejected graph:
\begin{verbatim}
variables:
  A:
    kind: record_set
    role: contract_records
    filters:
      supplier = GLEAM SERVICE CENTRE LTD
    depends_on: []

relations: []

return:
  operation: count
  input: A
\end{verbatim}

The serialized object has no buyer entity variable and no relation.  Its stale template field says
\texttt{abstain}, while its operative return field says \texttt{count}.  The runtime uses the
compiled return operation and therefore executes the proposal.  The historical compact trace does
not retain a route tag, but the exact compiler equality means that no initial Stage 2 language model
call is needed or attributed to this item.  The configured Stage 2 prompt is relevant only when the
intent program fast path cannot produce a valid graph.

\section{First execution and layered verdict}

The executor resolves GLEAM SERVICE CENTRE LTD to a canonical supplier and exhaustively retrieves
35 matching awards.  The program returns 35.  Grounding, schema validity, operation support,
execution, nonempty evidence, and output shape all pass.  The online runtime therefore has no basis
for initiating a repair.

The offline teacher controller then compares 35 with the hidden oracle 1,623.  It records
\texttt{verified=true}, \texttt{oracle\_match=false}, and \texttt{answer=35}, then exports the plan
as a hard negative.  The error is semantic rather than syntactic.  The program correctly executes
the wrong population.

\begin{table}[h]
\centering
\small
\begin{tabular}{p{0.24\linewidth}p{0.28\linewidth}p{0.38\linewidth}}
\toprule
Stage & Observation & Evidence status \\
\midrule
Question & Hidden oracle 1,623 & Stored benchmark row \\
Initial proposal & Count supplier records & Stored teacher trace \\
Initial execution & 35; online checks pass & Stored teacher trace \\
External audit & 35 differs from 1,623 & Stored trace and hard negative \\
Repair input & Validation failed; 35 submitted & Audited prompt reconstruction \\
Repaired execution & $35\rightarrow2\rightarrow1{,}623$ & Stored target and deterministic replay \\
\bottomrule
\end{tabular}
\caption{The complete observation sequence.}
\end{table}

\section{Oracle hidden repair}

Oracle agreement is used only by the offline controller to decide whether a repair attempt should
be collected.  The fixed graph repair system message is
\begin{quote}\small
You are the repair reflector inside a verifier-guided KGQA graph planner. Diagnose the failed graph
plan from the structured failure feedback, choose one allowed repair action, and return a repaired
understanding network plus repaired graph plan using the SAME schema as the planning step. Never
answer the question and never use hidden reference answers. Return strict JSON.
\end{quote}

For this item, the visible feedback is equivalent to
\begin{verbatim}
failure_stage: verifier
failure_reason: wrong_answer
submitted_answer: 35
grounding_issues: []
schema_errors: []
failed_checks: []
allowed_repair_actions:
  - fix_question_type
  - repair_constraints
  - swap_buyer_supplier
  - change_operator
  - change_operation
  - abstain
notes:
  - The submitted answer did not match the hidden verifier.
  - The reference/oracle answer is intentionally hidden from the reflector.
  - Repair the plan using only the question and trace summary.
\end{verbatim}

The full repair object also contains the failed graph and online check summaries.  Before prompt
construction, fields named \texttt{oracle\_answer}, \texttt{reference\_answer},
\texttt{gold\_answer}, and \texttt{expected\_answer}, as well as keys containing
\texttt{oracle}, are removed.  The compact historical export records the feedback as
\texttt{submitted answer failed external validation}.  It does not retain the raw provider repair
completion byte for byte.

\section{Accepted repair and deterministic replay}

The accepted repair inserts the missing bridge:
\begin{verbatim}
a1 = record_set(supplier = GLEAM SERVICE CENTRE LTD)
b1 = distinct buyer(a1)
c1 = record_set(buyer in b1)
return count(c1)
\end{verbatim}

The stored graph contains two explicit relations.  The first binds the buyer field from
\texttt{a1} into entity set \texttt{b1}.  The second binds those buyers into target record set
\texttt{c1}.  Deterministic replay exposes the following intermediate populations:

\begin{table}[h]
\centering
\small
\begin{tabular}{clrp{0.43\linewidth}}
\toprule
Variable & Kind & Size & Interpretation \\
\midrule
\texttt{a1} & Record set & 35 & Awards whose supplier is GLEAM \\
\texttt{b1} & Entity set & 2 & Distinct buyers derived from those awards \\
\texttt{c1} & Record set & 1,623 & All records published by the two buyers \\
\bottomrule
\end{tabular}
\caption{Intermediate populations in the repaired execution.}
\end{table}

The two buyer identities are London Borough of Haringey, Passenger Transport Services, and Royal
Borough of Greenwich, managed by GS Plus Ltd.  They contribute 1,370 and 253 records respectively.
Their sum is 1,623.  A representative source record is
\texttt{contract:ocds-h6vhtk-035eee:022611-2022-1}, titled ``Regular Services'' and carrying CPV
60170000.  The replay trace retains the identifiers and provenance for every intermediate record;
printing all 1,623 target rows is unnecessary for interpreting the transition.

The repaired execution passes the same online checks as the rejected execution and additionally
agrees with the hidden oracle.  The controller exports the repaired plan as supervised repair data
and exports the repaired and rejected plans as an oracle curated preference pair.

\section{What changed}

The repair did not improve entity grounding, fix invalid syntax, or turn a failing program into an
executable one.  Those properties already held.  It changed the reasoning action by inserting the
omitted supplier to buyer to record bridge.  The result changed from the number of GLEAM awards to
the number of records published by buyers that had awarded GLEAM a contract.

This distinction is the purpose of the environment.  It exposes the proposed program, every
intermediate population, the final result, and the authority of each verdict.  The environment does
not itself prove that an agent has improved.  It makes the change and its consequences observable
enough to attribute a successful repair.

\section{Artifact locations}

All historical objects are joined by \texttt{L2::bridge\_join\_0583\#L2a}.

\begin{tabular}{p{0.48\linewidth}p{0.44\linewidth}}
\toprule
Artifact & Role \\
\midrule
\path{data/qa/cicada_core_v4/train.jsonl} & Question and hidden oracle \\
\path{data/qa/teacher_full_v1/traces.jsonl} & Briefing, rejected graph, answer, verdict \\
\path{data/qa/teacher_full_v1/hard_negatives.jsonl} & Offline rejected example \\
\path{data/qa/teacher_full_v1/repair_sft.jsonl} & Visible failure signal and accepted target \\
\path{data/qa/teacher_full_v1/dpo_pairs.jsonl} & Accepted and rejected preference pair \\
\path{src/procurement_graph/reasoning/typed_planning.py} & Prompt and repair message builders \\
\path{scripts/eval_targeted_v2.py} & Hidden mismatch feedback construction \\
\path{scripts/run_teacher.py} & Offline routing and export conditions \\
\bottomrule
\end{tabular}

\end{document}
